\section{Diskussion und Fazit}

\subsection{Zusammenfassung der Befunde}

Die Analysen zeigen, dass Sentiment-Informationen in Pythia-410M bereits im Embedding-Raum strukturiert sind und für Ein-Token-Prompts mit bis zu 97~Prozent Genauigkeit linear getrennt werden können. Über die Transformer-Schichten hinweg baut sich die Sentiment-Repräsentation schrittweise auf und erreicht ihre stärkste Ausprägung in den oberen Schichten 17 bis 23. Die Attention-Analyse zeigt auf Einzelbeispielen eine starke Fokussierung auf sentimenttragende Wörter, dieser Effekt ist über den CAD-Datensatz jedoch insgesamt schwach und verteilt. Activation Patching liefert schließlich kausale Evidenz dafür, dass sentimenttragende Wörter einen entscheidenden Beitrag zur finalen Vorhersage leisten.

\subsection{Hypothesen-Verdict}

\textbf{H1 (Klare Stimmungswörter): Bestätigt.}

Positive und negative Wörter sind bereits früh im Modell unterscheidbar. Sentiment-Signale werden in frühen bis mittleren Schichten sichtbar und das Patchen des sentimenttragenden Wortes kann die ursprüngliche Vorhersage weitgehend wiederherstellen.

\textbf{H2 (Ironie und Sarkasmus): Teilweise bestätigt.}

Die Ergebnisse deuten darauf hin, dass Pythia-410M überwiegend der wörtlichen Polarität folgt. Für die tatsächlich gemeinte ironische Polarität konnten keine stabilen internen Repräsentationen nachgewiesen werden. Aufgrund des begrenzten Umfangs der Ironie-Analyse bleibt die Evidenz jedoch eingeschränkt.

\textbf{H3 (Lokalisierbarkeit): Teilweise bestätigt.}

Activation Patching identifiziert die Position des sentimenttragenden Wortes als stärksten kausalen Einflussfaktor. Die Sentiment-Verarbeitung erstreckt sich jedoch über mehrere Schichten und kann nicht auf eine einzelne Schicht oder einen einzelnen Attention Head reduziert werden.

\subsection{Korrelation versus Kausalität}

Ein zentrales Ergebnis dieser Arbeit ist die Unterscheidung zwischen korrelativen und kausalen Analysen. Logit Lens, Linear Probing und die Attention-Analyse identifizieren Schichten und Komponenten, in denen Sentiment-Informationen sichtbar oder dekodierbar sind. Die stärksten Effekte wurden dabei überwiegend in den mittleren und oberen Schichten beobachtet. Activation Patching zeigt jedoch, dass die für die Vorhersage relevante Information bereits an der Repräsentation des sentimenttragenden Tokens in frühen Schichten vorhanden sein kann. Hohe Attention-Gewichte oder eine starke Logit-Trennung sind daher wertvolle Indikatoren, erlauben jedoch keine direkten Aussagen über Kausalität. Erst die Kombination mehrerer Interpretierbarkeitsmethoden ermöglicht ein vollständigeres Bild der internen Verarbeitung.

\subsection{Fazit}

Die Ergebnisse dieser Arbeit zeigen, dass Pythia-410M Sentiment-Informationen bereits im Embedding-Raum repräsentiert und diese über die Transformer-Schichten hinweg schrittweise zu einer immer stärkeren Sentiment-Repräsentation ausbaut. Logit Lens und Linear Probing zeigen eine zunehmende Trennung positiver und negativer Repräsentationen, wobei die stärksten Effekte in den oberen Schichten auftreten. Die Attention-Analyse identifiziert einzelne Attention Heads mit Fokus auf sentimenttragende Wörter, deutet jedoch auf eine verteilte Verarbeitung hin. Activation Patching liefert schließlich kausale Evidenz dafür, dass sentimenttragende Wörter einen direkten Einfluss auf die finale Modellvorhersage besitzen und dass sich dieser Einfluss auf bestimmte Tokenpositionen und Schichtbereiche zurückführen lässt. Insgesamt zeigt die Arbeit, dass die Kombination von Logit Lens, Attention Analysis, Linear Probing und Activation Patching komplementäre Einblicke in die interne Verarbeitung von Sentiment in Transformer-Modellen ermöglicht und sowohl die Stärken als auch die Grenzen der einzelnen Interpretierbarkeitsmethoden sichtbar macht.